% Figure: orifice plate with the vena contracta downstream and the static
% pressure profile dipping at the contraction, marking the cavitation margin.
\begin{figure}[htbp]
\centering
\begin{tikzpicture}[
  wall/.style={draw=engblack, very thick},
  jet/.style={draw=engblue, thick},
  flow/.style={draw=engblue, -{Stealth}, thick},
  label/.style={font=\scriptsize},
]
% ==== pipe walls ====
\draw[wall] (-4,1.1) -- (4,1.1);
\draw[wall] (-4,-1.1) -- (4,-1.1);
% ==== orifice plate at x=0 ====
\draw[wall, fill=enggray!25] (-0.12,1.1) -- (0.12,1.1) -- (0.12,0.45) -- (-0.12,0.45) -- cycle;
\draw[wall, fill=enggray!25] (-0.12,-1.1) -- (0.12,-1.1) -- (0.12,-0.45) -- (-0.12,-0.45) -- cycle;
\node[label, anchor=south] at (0,1.15) {orifice plate};
% ==== streamlines through the hole, contracting to a vena contracta ====
\draw[jet] (-4,0.9) .. controls (-1.2,0.85) and (-0.3,0.45) .. (0.12,0.42)
  .. controls (0.7,0.36) and (1.0,0.30) .. (1.1,0.30)
  .. controls (1.6,0.32) and (2.6,0.55) .. (4,0.7);
\draw[jet] (-4,-0.9) .. controls (-1.2,-0.85) and (-0.3,-0.45) .. (0.12,-0.42)
  .. controls (0.7,-0.36) and (1.0,-0.30) .. (1.1,-0.30)
  .. controls (1.6,-0.32) and (2.6,-0.55) .. (4,-0.7);
\draw[flow] (-3.2,0) -- (-2.0,0);
% vena contracta marker
\draw[engred, {Stealth}-{Stealth}, thin] (1.1,0.30) -- (1.1,-0.30);
\node[label, engred, anchor=west] at (1.2,0) {vena contracta $A_c$};
\node[label, anchor=north] at (0,-1.15) {area $A_2$};
\node[label, anchor=north] at (-2.6,-1.15) {area $A_1$, $p_1$};
% ==== static-pressure mini-plot below ====
\begin{scope}[yshift=-3.6cm]
\draw[enggray, -{Stealth}] (-4,0) -- (4.2,0) node[label, anchor=west] {$x$};
\draw[enggray, -{Stealth}] (-4,-0.2) -- (-4,1.8) node[label, anchor=south] {$p$};
% pressure: high upstream, deep dip at vena contracta, partial recovery
\draw[engblue, very thick]
  (-4,1.5) .. controls (-1.5,1.45) and (-0.5,1.0) .. (1.1,0.15)
  .. controls (2.0,0.6) and (3.0,0.95) .. (4,1.0);
\node[label, engblue, anchor=south west] at (-3.9,1.5) {$p_1$};
\node[label, engblue, anchor=north] at (1.1,0.1) {$p_{\min}$};
% vapour pressure margin line
\draw[engred, dashed, very thick] (-4,0.32) -- (4,0.32);
\node[label, engred, anchor=west] at (3.0,0.32) {$p_v$};
\draw[engamber, {Stealth}-{Stealth}, thin] (1.1,0.15) -- (1.1,0.32);
\node[label, engamber, anchor=west] at (1.2,0.24) {margin};
\end{scope}
\end{tikzpicture}
\caption[Orifice-plate flow and the cavitation margin]{Top: flow
through a thin orifice plate. The streamlines continue to contract
past the plate to a minimum-area \emph{vena contracta} of area
$A_c<A_2$, where the speed is highest and the static pressure lowest;
the discharge coefficient $C_d\approx0.6$ folds the contraction and a
small viscous loss into one factor. Bottom: the static pressure along
the centreline, high upstream, dipping to $p_{\min}$ at the vena
contracta, and recovering only partially downstream. The cavitation
margin is the gap between $p_{\min}$ and the vapour pressure $p_v$; a
metering orifice sized for a large pressure drop can push $p_{\min}$
below $p_v$ and cavitate, so the design discipline is to check the
margin at the worked flow. Computed from continuity, Bernoulli, and
the vena-contracta contraction.}
\label{fig:vol03ch11:orifice-vena-contracta}
\end{figure}
